\documentclass[UTF8,11pt]{ctexart}
\usepackage[a4paper,left=2.55cm,right=2.55cm,top=2.5cm,bottom=2.55cm]{geometry}
\usepackage{amsmath,amssymb,mathtools}
\usepackage{array,booktabs,tabularx,longtable}
\usepackage{enumitem}
\usepackage{tikz}
\usetikzlibrary{arrows.meta,positioning,calc}
\usepackage{xcolor,tcolorbox}
\tcbuselibrary{breakable,skins}
\usepackage{hyperref}
\usepackage{fancyhdr}
\usepackage{microtype}
\definecolor{deepblue}{HTML}{244B74}
\definecolor{paleblue}{HTML}{F1F6FC}
\definecolor{grayblue}{HTML}{F4F6F9}
\hypersetup{colorlinks=true,linkcolor=deepblue,urlcolor=deepblue,bookmarksnumbered=true}
\setlength{\parindent}{2em}
\setlength{\parskip}{.35em}
\setlist[itemize]{leftmargin=2em,itemsep=2pt,topsep=4pt}
\setlist[enumerate]{leftmargin=2em,itemsep=2pt,topsep=4pt}
\newtcolorbox{proofbox}[1][]{enhanced,breakable,colback=paleblue,colframe=deepblue!40,boxrule=.5pt,arc=1.5mm,left=3mm,right=3mm,top=2mm,bottom=2mm,#1}
\pagestyle{fancy}\fancyhf{}\fancyhead[L]{自适应分组测试：核心证明逻辑}\fancyhead[R]{\nouppercase{\leftmark}}\fancyfoot[C]{\thepage}\setlength{\headheight}{15pt}
\title{自适应分组测试论文核心证明思路\\[.5em]\large 从证明目标出发，逐步理解整条证明逻辑}
\author{}
\date{}
\begin{document}
\maketitle
\tableofcontents
\vspace{1em}
\noindent\textit{基于 \emph{When Pooling Cannot Save a Test: A Guided Proof for Adaptive Group Testing with a Known Defective Count}}\par

这次我们不按照论文原来的章节顺序，而是尝试站在研究者的角度，重新构建整条证明路线。

我们的原则是：

\textbf{先明确想证明什么，再考虑应该怎样证明；如果当前思路遇到了困难，就解释为什么需要引入新的工具，以及这个工具如何帮助我们继续向最终目标推进。}

因此，下面不仅要说明每个公式的来源，还要解释它为什么会在这个位置出现。

对于已经在信息论课程中学过的基础定理，我们不再重复完整证明；但对于连接整条证明的关键论证，尤其是结构递推、停止条件、概率模型比较、全树信息不等式及最终反证，我们会保留必要的推导。

\section{确立最终目标：为什么采用反证法？}
\subsection{我们究竟要证明什么？}
假设有 $n$ 件物品，其中恰好 $d$ 件是缺陷品，并且我们事先知道 $d$ 的准确值。

每次测试可以选择任意物品集合 $S$，得到一个二元回答：
\[
Y(S)=\begin{cases}0,&S\text{ 中没有缺陷品},\\1,&S\text{ 中至少有一件缺陷品}.\end{cases}
\]
允许下一次测试根据之前所有回答自适应地选择。

记 $M(d,n)$ 为所有正确的确定性自适应分组测试算法中，最小的最坏情况测试次数。

由于已知缺陷品总数，只需逐个测试前 $n-1$ 件物品，最后一件的状态就能自动确定。因此：
\[M(d,n)\le n-1.\]
论文要证明，当
\[n\ge2^{31},\qquad 0<d<n,\qquad\frac dn\ge c_\star,\]
其中
\[\boxed{c_\star=\frac{369056}{1000000}=0.369056}\]
时，实际上有：
\[\boxed{M(d,n)=n-1.}\]
也就是说，在这个范围内，任何自适应分组测试算法都无法保证最坏情况下比逐个测试节省至少一次测试。

\subsection{为什么反证法比较自然？}
我们已经知道：
\[M(d,n)\le n-1.\]
所以困难只在于证明：
\[M(d,n)\ge n-1.\]
一个自然的思路是：假设存在某个更好的算法，它在所有有效输入上最多使用 $n-2$ 次测试，然后证明这样的算法不可能存在。

为了让“节省测试次数”成为可以直接分析的数学量，定义：
\begin{equation}\boxed{\Delta(d,n)=n-1-M(d,n).}\tag{1}\end{equation}
这里：
\begin{itemize}
\item $\Delta(d,n)=0$：逐个测试已经最优。
\item $\Delta(d,n)\ge1$：存在算法在最坏情况下至少节省一次测试。
\end{itemize}
由于 $M(d,n)\le n-1$，节省量始终是非负整数。

于是，证明目标变成：
\[\boxed{\Delta(d,n)=0.}\]
我们假设结论不成立：
\begin{equation}\boxed{\Delta(d,n)\ge1.}\tag{2}\end{equation}
接下来，所有工作都围绕这个假设展开。

\subsection{直接分析原问题为什么困难？}
如果试图直接证明：所有密度大于 $0.369056$ 的问题都必须使用至少 $n-1$ 次测试，就需要分析各种可能的测试策略。

而测试池可以任意选择、任意重叠，后续查询又可以依赖全部历史。因此，直接对所有决策树进行结构分析非常困难。

作者采取了另一种策略：

\textbf{先利用分组测试本身的组合结构，把一个假设存在节省的高密度问题转化成规模更小、密度更适合概率分析的问题。}

然后再利用信息论证明：这个新问题不可能具有递推所强迫的节省量。

这就把证明拆成了两个可以配合的任务：
\begin{enumerate}
\item 将原问题转化成一个带有强节省量要求的新问题。
\item 证明这个新问题无法达到所要求的节省量。
\end{enumerate}
我们先解决第一个任务。

\section{第一个困难：怎样把高密度问题转化为更容易分析的问题？}
\subsection{为什么希望新问题的密度接近 $1/4$？}
这里先解释作者为什么特别关注 $1/4$。

假设每件物品独立地以概率 $p$ 成为缺陷品，并记：
\[q=1-p\]
为正常品概率。

那么，一个包含 $s$ 件物品的测试池全部为正常品的概率就是 $q^s$。

如果取：
\[p=\frac14,\qquad q=\frac34,\]
那么：
\[q^2=\frac9{16}>\frac12,\qquad q^3=\frac{27}{64}<\frac12.\]
所以在这种独立模型下：
\begin{itemize}
\item 两件物品组成的测试池，阴性概率大于 $1/2$；
\item 至少三件物品组成的测试池，阴性概率小于 $1/2$。
\end{itemize}
这意味着两类测试池的回答都具有一定的不平衡性，只是方向相反。

为什么这有用？

因为我们已经知道，测试结果越不平衡，其信息熵越小，信息缺口越大。

如果能够将这种不平衡性转化为真实模型下不可避免的损失，就有希望获得比普通熵下界更强的结论。

这里先把密度 $1/4$ 看作一个\textbf{适合进行信息论分析的目标区域}。后面我们会严格解释，为什么独立模型下的不平衡性可以用于真实的自适应测试问题。

现在首先要解决：原始问题的缺陷密度至少是 $0.369056$，怎样把它转移到密度约为 $1/4$ 的问题，同时保留原问题假设存在的测试节省？

\subsection{为什么需要研究问题之间的递推关系？}
我们希望找到一个关系：如果某个规模为 $(r,N)$ 的问题可以节省测试，那么某个规模更小的问题也必须能够节省测试。

这样，我们就可以不断缩小问题规模，并观察缺陷密度怎样变化。

为此，先考虑一个具有 $N$ 件物品、$r$ 件缺陷品的问题。

假设某个正确算法在根节点选择了第一次测试集合 $S$。

为了分析它的最坏情况测试次数，我们不需要同时研究所有回答分支。

只要找到其中一个足够困难的分支，就能够得到整个算法最坏测试次数的下界。

根据第一次测试池的大小，我们考虑两种情况。

\subsection{情况一：第一次只测试一件物品}
假设：
\[|S|=1.\]
如果回答阴性，那么这件物品被确定为正常品。

剩下的问题是：$N-1$ 件物品中，恰好 $r$ 件缺陷品。

因此，这个分支至少还需要 $M(r,N-1)$ 次测试。

再算上第一次测试：
\[\boxed{M(r,N)\ge1+M(r,N-1)}\]
对于以单件测试开始的算法成立。

这个推导比较容易理解。真正困难的是第二种情况。

\subsection{情况二：第一次测试至少两件物品}
假设：
\[|S|\ge2.\]
从 $S$ 中任选两件不同的物品，记为 $\xi,\eta$。

现在考虑第一次测试为阳性的分支，并把候选输入进一步限制为：
\[\text{$\xi,\eta$ 中至少有一件是缺陷品}.\]
因为 $\xi,\eta\in S$，所以这些输入一定产生阳性回答。

因此，这些输入构成阳性分支中的一个合法子问题。

但是，这个子问题与普通的“恰好 $r$ 件缺陷品”的完整问题不同。

原因是我们多知道了一个条件：
\[\xi,\eta\text{ 中至少有一件缺陷品}.\]
那么，它的测试复杂度是多少？

这就是当前遇到的第一个真正的数学障碍。

我们不能简单地删掉 $\xi$ 或 $\eta$，因为还不知道哪一件有缺陷。

为了解决这个困难，作者引入\textbf{两特殊标签恒等式}，即论文引理 10.3。

\subsection{两特殊标签恒等式解决了什么？}
设除了 $\xi,\eta$ 以外，还有 $m$ 件普通物品，组成集合 $U$。

考虑如下候选输入族：
\[
\begin{aligned}
\mathcal J={}&\{\{\xi\}\cup D:D\subseteq U,\ |D|=k\}\\
&\cup\{\{\eta\}\cup D:D\subseteq U,\ |D|=k\}\\
&\cup\{\{\xi,\eta\}\cup E:E\subseteq U,\ |E|=k-1\}.
\end{aligned}
\]
这三类输入分别表示：
\begin{itemize}
\item 只有 $\xi$ 是特殊标签中的缺陷品；
\item 只有 $\eta$ 是特殊标签中的缺陷品；
\item $\xi,\eta$ 都是缺陷品。
\end{itemize}
所以 $\mathcal J$ 恰好描述：总共有 $k+1$ 件缺陷品，并且 $\xi,\eta$ 中至少有一件缺陷品。

记 $\operatorname{comp}(\mathcal J)$ 为已知候选输入属于 $\mathcal J$ 时，使用合法 OR 测试精确识别输入所需的最优最坏情况测试次数。

作者证明：
\begin{equation}\boxed{\operatorname{comp}(\mathcal J)=1+M(k,m+1).}\tag{3}\end{equation}
这个结果的重要性在于：

\textbf{一个带有“两件指定物品至少一件缺陷”条件的问题，可以精确地转化为一个规模减少、缺陷数量也减少的标准分组测试问题，再加一次测试。}

为什么是严格等式，而不是仅仅一个上界？

容易的一边是先测试 $\xi$，从而得到：
\[\operatorname{comp}(\mathcal J)\le1+M(k,m+1).\]
困难的是反方向。

作者通过决策树模拟证明：从任意正确识别 $\mathcal J$ 的决策树，都能构造出一棵解决标准 $(k,m+1)$ 问题、且最坏深度至少减少 1 的合法决策树。

主要思想是：将标准问题的每个输入映射成原特殊标签问题中的一个代表输入。

普通物品测试的回答在两个输入之间保持一致。

当原树第一次查询特殊标签时，选择适当的代表映射，使这次查询的回答在所有代表输入上都确定为阳性，因此能够跳过该次测试。

后续涉及特殊标签的查询也通过固定回答或合法的标签替换进行模拟。

还要处理一种特殊情况：某条路径可能从头到尾都没有查询 $\xi,\eta$。

作者证明这种路径不可能恰好达到原决策树的最大深度。

原因是：普通查询无法区分仅在 $\xi,\eta$ 上不同的两类候选输入；而如果只剩下两特殊标签同时缺陷的候选类型，论文又通过候选集合的扩张性质证明，原树不可能在最深层才完成这类识别。

这样就保证了模拟后的每条路径都至少减少一次测试，而不只是某些路径减少一次。

因此：
\[\operatorname{comp}(\mathcal J)\ge1+M(k,m+1),\]
两边合起来才得到式 (3)。

引理的完整决策树模拟细节可以在原文第 10.3 节查看。这里需要抓住的是：\textbf{作者没有假设最优算法一定先测试特殊标签，而是证明了对任意合法测试决策树都成立的结构恒等式。}

\subsection{由此得到结构递推}
现在回到原来的 $(r,N)$ 问题。

若第一次测试至少两件物品，则取：
\[m=N-2,\qquad k=r-1.\]
于是式 (3) 给出：阳性分支中选定的子问题至少需要 $1+M(r-1,N-1)$ 次测试。

再加上根节点的第一次测试：
\[M(r,N)\ge2+M(r-1,N-1).\]
结合第一次测试为单件物品的情况，就得到：
\begin{equation}
\boxed{M(r,N)\ge\min\{1+M(r,N-1),\,2+M(r-1,N-1)\}.}
\tag{4}
\end{equation}
它在 $r\ge2,\ N\ge r+2$ 时成立。

这就是论文命题 10.4。

这里取最小值，是因为我们不知道最优算法会采用哪一类首次测试。对两种情况分别建立下界后，能够统一保证的就是两者的较小值。

\textbf{至此，我们获得了连接不同规模分组测试问题的关键工具。}

\section{结构递推如何强迫更小的问题节省更多测试？}
\subsection{把递推改写成关于节省量的关系}
我们现在真正关心的是：
\[\Delta(r,N)=N-1-M(r,N).\]
将其代入式 (4)。

首先：
\[1+M(r,N-1)=N-1-\Delta(r,N-1).\]
其次：
\[2+M(r-1,N-1)=N-\Delta(r-1,N-1).\]
所以：
\begin{equation}
\boxed{\Delta(r,N)\le\max\{\Delta(r,N-1),\,\Delta(r-1,N-1)-1\}.}
\tag{5}
\end{equation}
假设：
\[\Delta(r,N)\ge s.\]
那么，右边至少有一项不小于 $s$。

因此，至少能够推出下面两个结论中的一个：
\[\Delta(r,N-1)\ge s,\]
或者：
\[\Delta(r-1,N-1)\ge s+1.\]
这意味着我们总能进行以下两种参数转移中的一种：
\begin{center}
\begin{tabularx}{\linewidth}{c>{\centering\arraybackslash}X>{\centering\arraybackslash}X}
\toprule 转移&参数变化&节省量的变化\\\midrule
A&$(r,N)\to(r,N-1)$&至少保持不变\\
B&$(r,N)\to(r-1,N-1)$&至少增加 1\\\bottomrule
\end{tabularx}
\end{center}
这里的“转移”只是利用结构递推，在不同的完整分组测试问题之间寻找满足某个下界的新问题。

\textbf{它不是一个实际测试算法，也不表示我们能够免费从未知配置中删除物品。}

\subsection{为什么不能无限递推？}
现在我们已经有了一种办法：从假设满足 $\Delta(d,n)\ge1$ 的原问题开始，不断减小物品数量。

但随之出现一个新困难：式 (5) 并不是在所有参数下都成立。

它要求：
\[r\ge2,\qquad N\ge r+2.\]
因此，递推可能在 $r=1$ 或 $N=r+1$ 时失效。

如果在密度降到 $1/4$ 之前就遇到这些边界情况，整条证明路线就会中断。

这就是为什么作者需要进行\textbf{边界分析}。

\subsection{边界分析：证明递推不会提前中断}
论文首先证明：
\[M(N-1,N)=N-1,\]
以及：
\[M(1,2)=1,\qquad M(1,3)=2.\]
这些等式的意义是：当只有一件正常品，或者问题退化为上述两个最简单的一缺陷品问题时，逐个测试已经达到最坏情况最优。

因此，它们全部满足：
\[\Delta=0.\]
但是，我们的递推过程从 $\Delta(d,n)\ge1$ 开始。

每一次转移都保证节省量下界不减小。

所以，在整个过程中：
\[\Delta(r,N)\ge1.\]
这意味着，不可能到达任何满足 $\Delta=0$ 的边界问题。

特别地：如果 $r=1$ 且 $r/N>1/4$，那么只能有 $N=2$ 或 $N=3$。

这两个问题的节省量都是零，不可能成为递推的终点。

如果：
\[N=r+1,\]
则只剩下一件正常品，也有 $\Delta=0$。

同样不可能到达。

另一方面，每次转移都使 $N$ 减少 1，因此不可能无限持续。

所以：

\textbf{只要密度还大于 $1/4$，递推就不会因为碰到不适用的边界而停止；而递推又不可能无限进行，因此一定会到达密度首次不超过 $1/4$ 的问题。}

这就是边界分析解决的问题。

\subsection{停止时的密度为什么足够接近 $1/4$？}
现在已经知道递推一定可以进行到密度不超过 $1/4$。

但还有一个困难：由于 $r,N$ 都是整数，密度可能从略大于 $1/4$ 一下跳到明显小于 $1/4$。

而我们希望最终的新问题密度仍然非常接近 $1/4$。

因此，需要进行\textbf{整数舍入控制}。

首先观察两种转移对密度的影响。

转移 A：
\[(r,N)\to(r,N-1)\]
会提高密度，因为：
\[\frac r{N-1}>\frac rN.\]
因此，第一次使密度降到 $1/4$ 以下的转移，只可能是 B：
\[(r+1,N+1)\to(r,N).\]
由于停止前密度仍然大于 $1/4$：
\[\frac{r+1}{N+1}>\frac14.\]
整理得到：
\[4(r+1)>N+1,\]
所以：
\[4r>N-3.\]
两边除以 $4N$：
\[\frac rN>\frac14-\frac3{4N}.\]
而停止条件告诉我们：
\[\frac rN\le\frac14.\]
因此：
\begin{equation}\boxed{\frac14-\frac3{4N}<\frac rN\le\frac14.}\tag{6}\end{equation}
这个结论说明：停止时密度与 $1/4$ 的差距最多只有 $3/(4N)$。

也就是说，只要 $N$ 足够大，停止密度就足够接近 $1/4$。

但这里又出现了一个问题：

\textbf{原始 $n$ 很大，并不自动意味着递推后的 $N$ 仍然很大。}

因此，必须进一步证明停止后还保留了足够多的物品。

\subsection{为什么停止后仍然有 $N>2^{29}$？}
假设递推过程中发生了：
\begin{itemize}
\item $j$ 次转移 B，即缺陷数和物品数同时减一；
\item $s$ 次转移 A，即只有物品数减一。
\end{itemize}
那么：
\[r=d-j,\qquad N=n-j-s.\]
每一次转移 B 都使节省量下界增加一次。

所以从最初的 $\Delta(d,n)\ge1$ 出发，得到：
\begin{equation}\boxed{\Delta(r,N)\ge1+j.}\tag{7}\end{equation}
另一方面：
\[\Delta(r,N)=N-1-M(r,N).\]
由于测试次数不可能为负：
\[M(r,N)\ge0,\]
因此：
\[\Delta(r,N)\le N-1.\]
结合：
\[1+j\le\Delta(r,N)\le N-1,\]
得到：
\[j\le N-2.\]
又因为停止时：
\[r\le\frac N4,\]
以及：
\[d=r+j,\]
所以：
\[d\le\frac N4+N-2=\frac{5N}{4}-2.\]
整理得到：
\[N\ge\frac45(d+2).\]
原始问题满足：
\[d\ge c_\star n,\]
因此：
\[N\ge\frac45(c_\star n+2).\]
而：
\[c_\star=0.369056>\frac5{16}.\]
于是：
\[\frac45c_\star>\frac14.\]
所以：
\[N>\frac n4.\]
由于：
\[n\ge2^{31},\]
最终得到：
\begin{equation}\boxed{N>2^{29}.}\tag{8}\end{equation}
这就解决了规模问题。

代入式 (6)，由于：
\[\frac3{4N}<0.001,\]
得到：
\begin{equation}\boxed{0.249<\frac rN\le0.25.}\tag{9}\end{equation}
至此，我们证明：\textbf{递推后的问题不仅存在，而且规模足够大，密度也确实落在我们希望研究的 $1/4$ 附近。}

\subsection{新问题究竟被迫节省多少次测试？}
现在还需要解决最后一个问题：

我们知道：
\[\Delta(r,N)\ge1+j,\]
但 $j$ 究竟有多大？

如果 $j$ 很小，这个下界可能没有什么用。

因此，我们希望证明：
\[j\ge gN\]
对于某个正常数 $g$ 成立。

这样就能够得到一个线性阶的节省量要求。

从参数变化关系：
\[r=d-j,\qquad N=n-j-s,\]
以及 $s\ge0$，得到：
\[N\le n-j=n-d+r.\]
定义原始密度：
\[c=\frac dn,\]
以及停止后的密度：
\[p=\frac rN.\]
于是：
\[N\le n(1-c)+pN.\]
整理：
\[\frac nN\ge\frac{1-p}{1-c}.\]
另一方面：
\[
\begin{aligned}
\frac jN&=\frac{d-r}{N}\\
&=c\frac nN-p\\
&\ge\frac{c(1-p)}{1-c}-p\\
&=\frac{c-p}{1-c}.
\end{aligned}
\]
由于：
\[c\ge c_\star,\qquad p\le\frac14,\]
而函数 $\frac{c-p}{1-c}$ 随 $c$ 增大而增大、随 $p$ 增大而减小，所以：
\[\frac jN\ge\frac{c_\star-1/4}{1-c_\star}.\]
现在，我们终于引入常数：
\begin{equation}\boxed{g_\star:=\frac{c_\star-1/4}{1-c_\star}=\frac{7441}{39434}.}\tag{10}\end{equation}
因此：
\[j\ge g_\star N.\]
再代入式 (7)：
\begin{equation}\boxed{\Delta(r,N)\ge1+g_\star N.}\tag{11}\end{equation}
这里：
\[g_\star\approx0.188695.\]
\textbf{这个常数不是人为选择的，也不是某个信息论不等式中的参数。}

它来自原始密度下界 $c_\star$、目标停止密度 $1/4$，以及结构递推过程中参数之间的关系。

它衡量的是：为了使密度从至少 $0.369056$ 降到至多 $1/4$，至少需要发生多少次降低缺陷数量的转移，相对于最终规模 $N$ 的比例。

而每一次这种转移都会使节省量下界增加一次。

因此，结构递推强迫新问题节省至少约 $0.188695N$ 次测试。

\section{第一条证明主线完成：我们现在真正需要证明什么？}
到这里，我们已经完成了结构性论证。

从假设：
\[\Delta(d,n)\ge1\]
出发，我们证明存在一个新问题 $(r,N)$，满足：
\[\boxed{N>2^{29},\qquad0.249<\frac rN\le0.25,}\]
并且：
\[\boxed{\Delta(r,N)\ge1+g_\star N.}\]
现在证明目标发生了变化。

我们不再需要直接分析原始高密度问题的所有算法。

我们只需要证明：

\textbf{对于密度接近 $1/4$ 且规模足够大的问题，任何正确分组测试算法都无法节省 $1+g_\star N$ 这么多次测试。}

这就是为什么作者接下来转向信息论。

注意，信息论现在有了一个非常明确的任务：

不是随便建立某个测试次数下界，而是要得到一个足够强的上界，限制：
\[\Delta(r,N)=N-1-M(r,N).\]
只要最终得到的节省量上界严格小于 $1+g_\star N$，就能够完成反证。

接下来，我们沿着这个目标继续推进。

\section{第二条主线开始：为什么普通熵下界不足以排除这个节省量？}
\subsection{先尝试最简单的信息论方法}
现在我们研究停止后的问题：
\[(r,N),\qquad p=\frac rN\approx\frac14.\]
为了利用概率论，我们从所有恰好包含 $r$ 件缺陷品的配置中均匀随机抽取一个。

记这种概率分布为 $P$，并令：
\[X=\text{随机缺陷品配置}.\]
因为有 $\binom Nr$ 种等可能的配置，所以初始熵为：
\begin{equation}\boxed{H=H_P(X)=\log_2\binom Nr.}\tag{12}\end{equation}
固定任意一个正确的确定性自适应测试算法，记：
\[T=\text{该算法实际执行的测试次数},\]
\[L=\mathbb E_P[T].\]
这里的 $L$ 是平均测试次数，而不是最坏情况测试次数。

不过：
\[L\le M(r,N)\]
对于最优最坏情况算法也成立。

因此，如果能证明所有正确算法都满足：
\[L\ge C,\]
那么最优最坏情况测试次数也必须满足：
\[M(r,N)\ge C.\]
这就说明：\textbf{虽然最终目标是最坏情况复杂度，但我们仍然可以通过对有效输入取平均来建立下界。}

\subsection{普通熵下界能给出什么？}
我们已经知道：
\[L\ge H.\]
因此：
\[M(r,N)\ge H.\]
而：
\[\Delta(r,N)=N-1-M(r,N).\]
于是：
\[\Delta(r,N)\le N-1-H.\]
令：
\[h(p)=-p\log_2p-(1-p)\log_2(1-p).\]
我们在第十一课证明过：
\[0\le Nh(p)-H\le\log_2(N+1).\]
因此：
\[H=Nh(p)-O(\log N).\]
在 $p\approx1/4$ 时，普通熵下界允许的节省量大约为：
\[\Delta(r,N)\lesssim N(1-h(1/4)).\]
具体计算：
\[1-h(1/4)\approx0.1887218755.\]
但是，结构递推给出的要求是：
\[\Delta(r,N)\ge1+g_\star N,\]
其中：
\[g_\star\approx0.1886950347.\]
比较两个线性系数：
\[
\underbrace{0.1887218755}_{\text{普通熵允许的主要节省量系数}}
>
\underbrace{0.1886950347}_{\text{结构递推强迫的节省量系数}}.
\]
这意味着：\textbf{仅靠普通熵下界，并不能排除结构递推所要求的节省量。}

两个结论仍然可能同时成立，所以还没有矛盾。

不过，两个系数的差距已经很小：
\[(1-h(1/4))-g_\star\approx2.6841\times10^{-5}.\]
这给我们提供了一个重要启发：

\textbf{我们不需要把熵下界改善很多，只要得到一个足够精确、能够略微限制最大节省量的加强下界，就可能使两个系数的大小关系发生逆转。}

这就是引入额外信息缺口的真正原因。

\subsection{为什么研究信息缺口？}
我们定义：
\begin{equation}\boxed{\Gamma=L-H.}\tag{13}\end{equation}
因为 $L\ge H$，所以 $\Gamma\ge0$。

在第九课，我们已经证明：
\begin{equation}\boxed{\Gamma=\sum_vP(v)(1-h(\pi_v)).}\tag{14}\end{equation}
这里：
\begin{itemize}
\item $v$ 遍历决策树的内部测试节点；
\item $P(v)$ 表示真实模型 $P$ 下到达节点 $v$ 的概率；
\item $\pi_v=P(\text{阴性}\mid\text{到达 }v)$；
\item $1-h(\pi_v)$ 表示该节点的信息缺口。
\end{itemize}
因此：
\[\Gamma\]
恰好是整棵决策树的信息缺口按照真实访问概率加权后的总和。

如果 $\Gamma=0$，说明所有正访问概率的测试节点都必须完全平衡。

如果 $\Gamma>0$，就意味着算法的平均测试次数严格超过普通熵下界。

所以我们现在的任务变成：
\[\boxed{\text{证明任何合法分组测试决策树，都必须产生足够大的 }\Gamma.}\]
但是，怎样证明？

直接分析每个节点的条件阴性概率非常困难。

因为算法是自适应的，某个测试节点可能在复杂的历史之后才被访问。

此时剩余候选配置具有复杂的条件分布。

因此，我们需要改变分析方式。

\section{怎样把“合法分组测试的结构限制”转化成可以计算的数量？}
\subsection{先规范化测试池}
这里我们先不引入任何新的概率模型。

因为当前最重要的问题是：

\textbf{究竟应该从决策树的什么结构入手，才能证明额外信息缺口不可避免？}

固定一棵有限、正确的确定性测试决策树。

对于内部节点 $v$，记：
\[S_v=\text{该节点原本测试的物品集合}.\]
假设沿着到达 $v$ 的历史路径，此前某些测试给出了阴性回答。

在真实无噪声模型下，任何一个此前被阴性测试池包含的物品，都已经确定为正常品。

记这些物品的并集为 $Z_v$。

那么，我们可以把当前测试池替换成：
\begin{equation}\boxed{\widehat S_v=S_v\setminus Z_v.}\tag{15}\end{equation}
因为删掉的物品已经确定正常，不可能影响当前 OR 测试的回答。

论文证明，这种规范化对于任何真实输入都不会改变算法的执行路径和测试次数。

因此，我们以后只研究规范化测试池 $\widehat S_v$。

现在将它们分为：
\[
\begin{cases}
\text{小测试池},&|\widehat S_v|\le2,\\
\text{大测试池},&|\widehat S_v|\ge3.
\end{cases}
\]
注意，空测试池也归入小测试池。

\subsection{为什么按大小分类还不够？}
对于大小为 1 或 2 的测试池，我们还可以进一步区分：它是否包含以前已经在小测试池中出现过的物品？

例如，假设此前的小测试池已经包含物品 1、2。

现在测试：
\[\{2,3\}.\]
那么只有物品 3 是第一次出现在小测试池中。

而如果当前测试：
\[\{4,5\},\]
并且物品 4、5 此前都没有出现在小测试池中，那么这次测试同时引入了两件新物品。

这两种情况在计数时存在重要区别。

因此，作者特别定义：

\textbf{全新物品对（fresh pair）}：当前规范化测试池恰好有两件物品，而且这两件物品此前都没有在任何小规范化测试池中出现过。

它们以前可能出现在大测试池中。

这个区别在后面构造人工参考实验时非常重要。

\subsection{定义三个需要控制的数量}
现在考虑任意一条真实执行路径。

我们记录：
\[a=\text{这条路径上的大测试池数量},\]
\[b=\text{这条路径上的全新物品对数量},\]
\[w=\text{从未出现在小测试池中的物品数量}.\]
因为输入配置是随机的，所以 $a,b,w$ 也可能随输入变化。

我们在真实模型 $P$ 下求期望：
\begin{equation}\boxed{A=\mathbb E_P[a],\quad B=\mathbb E_P[b],\quad W=\mathbb E_P[w].}\tag{16}\end{equation}
因此：
\begin{itemize}
\item $A$ 是大测试池的平均执行次数；
\item $B$ 是全新物品对测试的平均执行次数；
\item $W$ 是真实执行结束时，平均还有多少物品从未出现在小测试池中。
\end{itemize}
为什么突然关注这三个量？

因为作者希望通过它们，直接约束总测试次数。

接下来我们先证明一个不需要任何信息论工具的组合不等式。

\subsection{只考虑一条路径：怎样计算物品第一次出现在小测试池中的次数？}
固定一条真实执行路径。

假设这条路径总共执行了 $t$ 次测试，其中 $a$ 次是大测试池。

那么小测试池共有 $t-a$ 次。

现在，我们统计不同物品第一次出现在小测试池中的情况。

对于每次小测试：
\begin{itemize}
\item 空测试池不会引入新物品；
\item 单件测试最多引入一件新物品；
\item 如果一个两件物品测试池不是全新物品对，那么最多引入一件新物品；
\item 如果是全新物品对，那么能够同时引入两件新物品。
\end{itemize}
因此，我们可以先给每次小测试池一个最多引入一件新物品的额度，再给每个全新物品对额外增加一个额度。

所以，这条路径上所有小测试池覆盖的不同物品总数，最多为：
\[t-a+b.\]
另一方面，$w$ 件物品从未出现在小测试池中。

因此，真正被小测试池覆盖的物品数恰好为：
\[N-w.\]
所以：
\[N-w\le t-a+b.\]
整理：
\begin{equation}\boxed{N-t\le b+w-a.}\tag{17}\end{equation}
这个不等式非常重要。

因为它对每条真实执行路径都成立，而不依赖于具体缺陷品配置。

现在对模型 $P$ 求期望：
\begin{equation}\boxed{N-L\le B+W-A.}\tag{18}\end{equation}
这里我们终于找到一个很好的突破口。

原来我们想约束 $L=\text{平均测试次数}$。

现在发现，只要能够控制 $B+W-A$，就能够约束 $L$。

\textbf{因此，证明任务进一步转化为：如何利用信息论证明 $B+W-A$ 不可能太大？}

\subsection{这个新任务有什么好处？}
因为 $A,B,W$ 具有明确的物理意义。

我们不再需要直接分析所有可能的决策树形状，而是可以尝试：
\begin{itemize}
\item 控制算法平均使用多少次大测试池；
\item 控制算法平均使用多少次全新物品对；
\item 控制算法平均留下多少从未出现在小测试池中的物品。
\end{itemize}
接下来，如果能够把这些数量与信息缺口 $\Gamma$ 联系起来，就能够通过式 (18) 得到我们需要的平均测试次数下界。

不过，这三个量不能用同一种简单方法控制。

我们先从大测试池开始。

\section{怎样控制大测试池？为什么需要独立伯努利模型？}
\subsection{真实模型下计算条件概率的困难}
在真实模型 $P$ 下，缺陷品数量始终等于 $r$。

因此，各物品状态并不独立。

此外，如果已经观察到许多测试回答，那么当前剩余候选配置还会受到更多条件限制。

所以，即使我们知道某个测试池包含至少三件物品，也不能简单写：
\[P(\text{阴性}\mid\text{历史})=(1-p)^3.\]
因为这个等式要求相关物品在条件分布下仍然独立，而这种独立性通常不成立。

我们已经确定：大测试池在某个独立模型下容易分析。

所以，自然想到：

\textbf{先在一个容易计算的独立模型中研究大测试池，再寻找数学工具，把得到的结果传递回真实模型。}

这就是引入第二个概率模型的原因。

\subsection{引入独立伯努利模型 $Q$}
现在令：
\[p=\frac rN,\qquad q=1-p.\]
定义两个输入概率模型：

真实模型 $P_X$：
\[
P_X(D)=\begin{cases}\displaystyle\frac1{\binom Nr},&|D|=r,\\0,&|D|\ne r.\end{cases}
\]
独立模型 $Q_X$：
\begin{equation}\boxed{Q_X(D)=p^{|D|}q^{N-|D|}.}\tag{19}\end{equation}
也就是说，在 $Q$ 下，每件物品独立地以概率 $p$ 成为缺陷品。

但这里要明确：

\textbf{我们没有改变原始分组测试问题的输入分布，也没有让真实算法改为处理独立模型。}

$Q$ 只是一个辅助分析模型。

我们固定同一棵决策树，在 $P$ 和 $Q$ 下分别运行真实 OR 测试。

为保证两个模型下的执行都具有定义，对于原来在恰好 $r$ 件缺陷品的承诺下不可能出现的回答分支，可以添加直接终止的辅助叶节点。

这不会改变任何真实有效输入下的执行路径。

于是，同一棵树产生两个完整路径分布：
\[P_Y,\qquad Q_Y.\]
其中 $Y$ 表示完整的终止回答序列。

\subsection{两个模型不同，为什么还能比较？}
我们在第十一课已经证明：
\[P_X=Q_X(\cdot\mid |X|=r).\]
因此：
\[D_{\mathrm{KL}}(P_X\|Q_X)=-\log_2Q_X(|X|=r).\]
记：
\begin{equation}\boxed{K:=D_{\mathrm{KL}}(P_X\|Q_X).}\tag{20}\end{equation}
由于：
\[Q_X(|X|=r)=\binom Nr p^rq^{N-r},\]
所以：
\[
\begin{aligned}
K&=-\log_2\binom Nr-r\log_2p-(N-r)\log_2q\\
&=Nh(p)-H.
\end{aligned}
\]
因此：
\begin{equation}\boxed{K=Nh(p)-H\le\log_2(N+1).}\tag{21}\end{equation}
这里再次强调：

$K$ 是两个输入概率模型之间的 KL 散度，不是测试次数。

它在后面起到一种“累计模型差异预算”的作用。

为什么这么说？

因为对于同一棵确定性无噪声决策树，完整路径 $Y$ 是输入 $X$ 的确定性函数。

所以，根据数据处理不等式：
\[D_{\mathrm{KL}}(P_Y\|Q_Y)\le D_{\mathrm{KL}}(P_X\|Q_X)=K.\]
再利用决策树版本的 KL 散度链式法则：
\begin{equation}\boxed{\sum_vP(v)d(\pi_v\|\rho_v^Q)\le K.}\tag{22}\end{equation}
其中：
\[\pi_v=P(\text{阴性}\mid v),\]
\[\rho_v^Q=Q(\text{阴性}\mid v),\]
而：
\[d(u\|v)=u\log_2\frac uv+(1-u)\log_2\frac{1-u}{1-v}\]
是两个二元概率分布之间的 KL 散度。

这里所有节点访问概率都在真实模型 $P$ 下计算。

这个公式的意义是：

\textbf{不管决策树有多少节点，沿整棵树按照真实访问概率加权后的模型差异总和，都不能超过同一个 $K$。}

这不是说每个节点都能单独使用一次 $K$ 的预算。

整棵树中的所有节点共享这一总额。

现在，我们拥有了控制两个模型差异的工具。

但还缺少另外一件事：

我们要证明的是信息缺口 $\Gamma$，而式 (22) 控制的是 KL 散度。

这两个量之间怎样联系？

\subsection{为什么需要加权二元不等式？}
假设在某个大测试池节点，比较模型 $Q$ 的阴性概率 $\rho_v^Q$ 明显小于 $1/2$。

我们希望推出真实模型下该节点也具有信息缺口。

但这并不能直接成立。

因为真实模型可能恰好满足：
\[\pi_v=\frac12,\]
使得：
\[1-h(\pi_v)=0.\]
不过，如果 $\pi_v=1/2$ 而 $\rho_v^Q$ 远离 $1/2$，那么两个模型之间必然存在正的 KL 散度。

反过来，如果两个模型在该节点的回答概率十分接近，那么真实回答概率也必须远离 $1/2$，从而产生信息缺口。

因此，产生一种自然的互补关系：

\textbf{比较模型明显不平衡时，真实算法要么付出信息缺口，要么付出模型差异。}

为了严格描述这种互补关系，论文引入加权二元不等式：
\begin{equation}
\boxed{\omega(1-h(\pi))+d(\pi\|\rho)\ge\frac{2\omega}{(\omega+1)\ln2}\left(\rho-\frac12\right)^2.}
\tag{23}
\end{equation}
其中 $\omega>0$ 是可以选择的权重。

这就是论文引理 6.3。

我们在第十二课已经学习了它的证明，这里不再重复。

只需要记住：如果能够证明比较模型的阴性概率满足：
\[\left|\rho-\frac12\right|\ge\delta>0,\]
那么：
\[\omega(1-h(\pi))+d(\pi\|\rho)\ge\frac{2\omega\delta^2}{(\omega+1)\ln2}.\]
因此，左边两种代价不能同时消失。

\textbf{它正好解决了当前的问题：把比较模型中的回答不平衡，转化成真实模型的信息缺口或 KL 散度代价。}

\subsection{但自适应历史会不会破坏 $Q$ 的独立性？}
现在还剩下一步。

我们需要证明大测试池在模型 $Q$ 下的条件阴性概率确实远离 $1/2$。

如果没有任何测试历史，那么：
\[Q(\text{测试池阴性})=q^{|\widehat S_v|}.\]
可是到达某个节点时，我们已经观察了此前的回答。

虽然输入在最初是独立的，但条件于阳性回答后，各物品状态不一定仍然独立。

因此，不能直接断言：
\[\rho_v^Q=q^{|\widehat S_v|}.\]
这是当前遇到的又一个困难。

为了解决它，作者利用独立伯努利分布的\textbf{单调事件关联不等式}。

其关键结论是：如果 $I$ 是一个递增事件，$J$ 是一个递减事件，而且 $Q(I)>0$，则：
\begin{equation}\boxed{Q(J\mid I)\le Q(J).}\tag{24}\end{equation}
在这里：
\begin{itemize}
\item 递增事件指增加缺陷品不会使事件由真变假；
\item 递减事件指增加缺陷品不会使事件由假变真。
\end{itemize}
例如：“某个测试池至少包含一件缺陷品”是递增事件。

“某个测试池全部为正常品”是递减事件。

因此，该不等式表达了：在独立伯努利模型下，条件于某些递增的阳性约束，不会提高另一个测试池全部为正常品的概率。

现在回到某个大测试节点 $v$。

此前所有阴性测试确认正常的物品，已经组成集合 $Z_v$。

在模型 $Q$ 下，条件于这些物品确实正常后，剩余物品仍然相互独立。

而此前所有阳性测试，则对剩余物品施加了递增事件约束。

因此，利用式 (24)，可以证明：
\begin{equation}\boxed{\rho_v^Q\le q^{|\widehat S_v|}.}\tag{25}\end{equation}
注意，这里得到的是不等式，不是错误地假设条件独立后得到的等式。

\subsection{终于可以控制大测试池}
对于大测试池：
\[|\widehat S_v|\ge3.\]
而停止后的密度满足：
\[0.249<p\le0.25.\]
论文为了统一处理边界，使用稍宽的区间：
\[0.249\le p\le0.251.\]
于是：
\[q\le0.751.\]
因此：
\[\rho_v^Q\le q^{|\widehat S_v|}\le(0.751)^3<\frac{17}{40}.\]
所以：
\begin{equation}\boxed{\frac12-\rho_v^Q>\frac3{40}.}\tag{26}\end{equation}
这意味着，大测试池在独立模型 $Q$ 下的阴性概率确实与 $1/2$ 保持固定距离。

现在使用加权二元不等式，并取：
\[\omega=\frac1{10}.\]
论文通过数值估计得到，对每个大测试节点：
\[\frac1{685}<\frac1{10}(1-h(\pi_v))+d(\pi_v\|\rho_v^Q).\]
将其乘以真实模型下的节点访问概率 $P(v)$，再对所有大测试节点求和。

记：
\[\Gamma_{\mathrm{lg}}=\sum_{v\text{ 为大测试池}}P(v)(1-h(\pi_v)).\]
根据 $A$ 的定义和式 (22)：
\[\frac A{685}\le\frac{\Gamma_{\mathrm{lg}}}{10}+K.\]
整理：
\begin{equation}\boxed{A\le\frac{137}{2}\Gamma_{\mathrm{lg}}+685K.}\tag{27}\end{equation}
这个结论终于完成了我们当前的任务。

\textbf{如果算法平均执行很多次大测试池，那么就必须产生足够多的大测试池信息缺口，或者消耗大量模型 KL 散度。}

而累计 KL 散度最多为：
\[K\le\log_2(N+1).\]
所以，大测试池的使用次数已经受到约束。

不过，我们还没有完成整体信息论证明。

因为式 (18) 需要控制 $B+W-A$，而不是单独控制 $A$。

因此，我们必须继续处理小测试池。

\section{大测试池已经受控，为什么小测试池仍然是困难？}
\subsection{为什么不能对所有小测试池重复刚才的方法？}
我们已经知道，大小至少为 3 的测试池在模型 $Q$ 下具有明显的不平衡性。

那么，大小为 2 的测试池呢？

在没有任何历史条件时：
\[Q(\text{两件物品均正常})=q^2.\]
当 $p\approx1/4$ 时：
\[q^2\approx\frac9{16}>\frac12.\]
所以，它同样是不平衡的。

乍看之下，我们似乎可以直接重复刚才的分析。

但这里存在一个问题：

\textbf{在模型 $Q$ 下，经过复杂的自适应历史后，当前两件物品不一定仍然相互独立。}

例如，之前某个大测试池已经得到阳性回答。

如果当前这两件物品曾经包含在那个大测试池中，那么它们的条件概率可能已经受到限制。

因此，即使这两件物品以前没有单独测试过，也不能直接写：
\[Q(\text{当前两件物品都正常}\mid\text{历史})=q^2.\]
对大测试池，我们只需要一个条件阴性概率的上界，可以通过单调事件不等式获得。

但对两件物品组成的测试池，我们希望得到 $\rho>1/2$ 的下界。

此前的阳性条件不仅不能帮助我们得到这个下界，反而可能使阴性概率减小。

这说明：

\textbf{模型 $Q$ 虽然适合控制大测试池，但不足以直接控制全新物品对。}

因此，需要设计另外一种辅助模型。

\subsection{我们希望新模型具有什么性质？}
回顾刚才遇到的困难。

小测试池的问题不是输入最初不独立。

模型 $Q$ 本来就是独立伯努利模型。

真正的问题是：此前的大测试池阳性回答可能在条件分布中引入依赖关系。

因此，我们希望构造一种概率实验，使得：
\begin{enumerate}
\item 输入物品状态仍然独立；
\item 小测试池继续使用真实 OR 回答；
\item 大测试池的回答不再对输入状态施加约束。
\end{enumerate}
如果能实现这三点，那么对于那些从未出现在小测试池中的物品，即使它们以前被大测试池包含过，其状态仍然可以保持独立。

这就自然引出了论文的第三种概率实验。

\subsection{构造人工参考实验 $R$}
在模型 $R$ 下：

首先，按照独立伯努利模型 $Q_X$ 抽取输入配置：
\[X_i\overset{\mathrm{i.i.d.}}{\sim}\operatorname{Bernoulli}(p).\]
然后，沿着已经固定的规范化决策树运行：

如果当前是小测试池：
\[|\widehat S_v|\le2,\]
则使用真实 OR 回答：
\[Y_v=\mathbf1_{\{X\cap\widehat S_v\ne\varnothing\}}.\]
如果当前是大测试池：
\[|\widehat S_v|\ge3,\]
则不计算真实 OR 结果，而是使用一枚新的独立公平硬币：
\[R(Y_v=0\mid\text{当前历史})=R(Y_v=1\mid\text{当前历史})=\frac12.\]
随后按照这个回答进入决策树的下一节点。

因此，三个模型现在具有明确的分工：
\begin{center}\small
\begin{tabularx}{\linewidth}{cXXX}\toprule
概率实验&输入分布&回答方式&研究目的\\\midrule
$P$&恰好 $r$ 件缺陷品，均匀分布&全部真实 OR&衡量真实算法的平均测试次数\\
$Q$&独立 Bernoulli($p$)&全部真实 OR&控制大测试池\\
$R$&独立 Bernoulli($p$)&小池真实，大池公平硬币&控制全新物品对及未被小测试触及的物品\\\bottomrule
\end{tabularx}
\end{center}
这里有两个特别重要的区别。

第一，$Q$ 和 $R$ 的输入边缘分布相同：
\[Q_X=R_X.\]
但它们产生测试回答的规则不同，因此路径分布通常不同：
\[Q_Y\ne R_Y.\]
第二，$R$ 并不是一个合法的无噪声分组测试算法。

因为它可能对实际包含缺陷品的大测试池给出阴性回答。

但是，这并不影响它作为比较概率模型使用。

\textbf{我们并没有要求真实算法使用人工回答；真实算法始终在模型 $P$ 下运行。}

\subsection{为什么 $R$ 不会破坏后续分析？}
这里需要处理一个容易被忽略的细节。

如果某个大测试池在 $R$ 下人工得到阴性回答，那么按照原决策树的历史规则，它的物品仍然会被加入集合 $Z_v$，并在后续规范化查询中被删除。

但是，这个人工阴性回答并不能证明那些物品真的正常。

因此，后续小测试池虽然仍然计算规范化集合的真实 OR 值，但不能把人工阴性当成关于输入 $X$ 的真实约束。

换句话说：

\textbf{规范化操作在 $R$ 中是沿固定决策树进行的形式操作，不意味着人工阴性回答真的确定了物品状态。}

这也是为什么必须明确构造人工参考实验，而不能简单地把它当成另一个合法测试算法。

\subsection{人工模型怎样保留条件独立性？}
固定某条从根节点到当前节点的回答历史：
\[\tau.\]
记：
\[a(\tau)=\text{历史中大测试池的数量}.\]
再定义一个事件：
\[E_\tau=\{\text{输入 }X\text{ 满足历史中所有小测试池的指定回答}\}.\]
注意，$E_\tau$ 只涉及小测试池的真实回答，不包含大测试池的人工回答所施加的输入约束。

对于任意具体输入配置 $D$，在模型 $R$ 下：
\begin{equation}
\boxed{R(X=D,\text{到达 }\tau)=Q_X(D)\,2^{-a(\tau)}\mathbf1_{E_\tau}(D).}
\tag{28}
\end{equation}
这个公式的三个因子分别对应：
\[\underbrace{Q_X(D)}_{\text{抽取输入 }D}\]
\[\underbrace{2^{-a(\tau)}}_{\text{大测试池公平硬币给出指定回答}}\]
和：
\[\underbrace{\mathbf1_{E_\tau}(D)}_{\text{小测试池回答必须与 }D\text{ 一致}}.\]
对 $D$ 求和：
\[R(\text{到达 }\tau)=2^{-a(\tau)}Q_X(E_\tau).\]
因此，只要该历史具有正概率：
\begin{equation}\boxed{R(X=D\mid\text{到达 }\tau)=Q_X(D\mid E_\tau).}\tag{29}\end{equation}
这个等式是模型 $R$ 的核心性质。

它说明：在固定历史条件下，输入的后验分布只受到此前小测试池回答的约束。

此前大测试池的人工回答没有提供任何额外的输入信息。

现在定义：
\[V_\tau=\text{历史中所有小规范化测试池的并集}.\]
事件 $E_\tau$ 只依赖于 $V_\tau$ 中的物品状态。

而模型 $Q_X$ 最初是独立伯努利分布。

因此，在条件于 $E_\tau$ 后，所有位于 $[N]\setminus V_\tau$ 中的物品仍然相互独立，并保持 Bernoulli($p$) 分布。

这就解决了我们之前遇到的条件独立性困难。

\subsection{终于可以准确计算全新物品对的阴性概率}
现在考虑某个全新物品对测试节点 $v$。

根据定义，当前测试池：
\[\widehat S_v=\{i,j\}\]
包含两件此前从未出现在小规范化测试池中的物品。

因此，在模型 $R$ 下，条件于到达 $v$：
\[X_i,X_j\]
仍然独立，而且各自以概率 $p$ 成为缺陷品。

所以当前测试阴性的条件概率恰好是：
\begin{equation}\boxed{\rho_v^R=R(\text{阴性}\mid v)=q^2.}\tag{30}\end{equation}
注意，这次是严格等式。

因为我们确实证明了相关物品在条件分布下仍然独立。

由于：
\[0.249\le p\le0.251,\]
所以：
\[q\ge0.749.\]
于是：
\[\rho_v^R=q^2\ge(0.749)^2>\frac{14}{25}.\]
因此：
\begin{equation}\boxed{\rho_v^R-\frac12>\frac3{50}.}\tag{31}\end{equation}
我们终于获得了想要的结论：

\textbf{全新物品对在人工模型 $R$ 下的阴性概率，与 $1/2$ 保持固定距离。}

现在又可以使用加权二元不等式了。

但是，与之前控制大测试池不同，这次还有一个新的困难：

我们并没有证明：
\[D_{\mathrm{KL}}(P_Y\|R_Y)\le K.\]
而且，这个不等式不能直接由数据处理不等式得到。

原因是 $P$ 与 $R$ 不再使用相同的回答生成规则。

因此，我们必须先研究：

\textbf{真实模型 $P$ 与人工模型 $R$ 之间的累计 KL 散度究竟是多少？}

\section{为什么必须研究联合分布，而不仅仅是路径分布？}
\subsection{为什么原来的数据处理方法不能直接使用？}
在模型 $P$ 与 $Q$ 下，固定输入 $X$ 后，测试回答完全相同。

因此，完整路径是两个模型共用的确定性函数：
\[Y=f(X).\]
这使我们能够使用数据处理不等式，得到：
\[D_{\mathrm{KL}}(P_Y\|Q_Y)\le K.\]
但对于人工模型 $R$，大测试池的回答由额外随机硬币决定。

因此，完整路径不再仅仅由输入 $X$ 决定。

所以，不能直接写：
\[D_{\mathrm{KL}}(P_Y\|R_Y)\le K.\]
我们必须重新计算模型差异。

这里作者想到：虽然 $P_Y$ 和 $R_Y$ 的概率比较复杂，但如果同时记录输入 $X$ 和完整路径 $Y$，那么联合概率反而容易计算。

这就是为什么要研究：
\[P_{X,Y},\qquad R_{X,Y}.\]

\subsection{先计算联合 KL 散度}
固定一个有效输入：
\[D\subseteq[N],\qquad |D|=r.\]
记：
\[y(D)=\text{真实算法在输入 }D\text{ 下的完整回答路径}.\]
再记：
\[a(D)=\text{这条真实路径上大测试池的数量}.\]
在真实模型 $P$ 下：
\[P_{X,Y}(D,y(D))=\frac1{\binom Nr}=2^{-H}.\]
而在人工模型 $R$ 下：首先抽到配置 $D$，其概率为：
\[Q_X(D)=p^rq^{N-r}=2^{-Nh(p)}.\]
其次，这条路径中每个小测试池的真实回答都与 $D$ 一致。

最后，每个大测试池必须通过公平硬币得到指定回答。

所以：
\[R_{X,Y}(D,y(D))=Q_X(D)\,2^{-a(D)}=2^{-Nh(p)-a(D)}.\]
因此：
\[
\begin{aligned}
\log_2\frac{P_{X,Y}(D,y(D))}{R_{X,Y}(D,y(D))}
&=\log_2\frac{2^{-H}}{2^{-Nh(p)-a(D)}}\\
&=Nh(p)-H+a(D)\\
&=K+a(D).
\end{aligned}
\]
在真实模型 $P$ 下对所有有效 $D$ 求期望：
\[
\begin{aligned}
D_{\mathrm{KL}}(P_{X,Y}\|R_{X,Y})
&=K+\mathbb E_P[a(X)]\\
&=K+A.
\end{aligned}
\]
所以：
\begin{equation}\boxed{D_{\mathrm{KL}}(P_{X,Y}\|R_{X,Y})=K+A.}\tag{32}\end{equation}
这是一条非常有用的精确恒等式。

它告诉我们：真实模型与人工模型的联合 KL 散度，恰好等于原来的输入模型差异 $K$，加上大测试池的平均数量 $A$。

为什么多出了 $A$？

因为每次人工大测试池回答使用一枚公平硬币。

对于真实路径中指定的大测试池回答，人工实验每次都需要付出一个概率 $1/2$ 的因子。

取负对数以后，每次恰好贡献 1 bit。

对真实路径求期望，就得到 $A$。

但请注意：式 (32) 描述的是\textbf{输入和完整路径的联合分布}。

我们真正需要用来控制全新物品对的，是路径分布之间的 KL 散度：
\[D_{\mathrm{KL}}(P_Y\|R_Y).\]
二者并不相同。

我们需要进一步将它们联系起来。

\subsection{为什么这时要使用 KL 散度链式法则？}
因为我们已经知道联合 KL 散度的精确值：
\[K+A.\]
而 KL 散度链式法则恰好能够把联合分布差异分解为：

第一部分：完整路径分布之间的差异；

第二部分：已经知道完整路径以后，两个模型对输入配置的条件分布差异。

因此：
\begin{equation}
\boxed{\begin{aligned}
K+A={}&D_{\mathrm{KL}}(P_Y\|R_Y)\\
&+\mathbb E_{Y\sim P_Y}D_{\mathrm{KL}}(P_{X\mid Y}\|R_{X\mid Y}).
\end{aligned}}
\tag{33}
\end{equation}
为了简化符号，定义：
\begin{equation}\boxed{D_R:=D_{\mathrm{KL}}(P_Y\|R_Y).}\tag{34}\end{equation}
因此：
\[K+A=D_R+\mathbb E_{Y\sim P_Y}D_{\mathrm{KL}}(P_{X\mid Y}\|R_{X\mid Y}).\]
现在出现一个新的可能性：

如果我们能够证明后面的条件 KL 散度至少为某个正数，那么就能得到 $D_R$ 的上界。

因为：
\[D_R=K+A-\mathbb E_{Y\sim P_Y}D_{\mathrm{KL}}(P_{X\mid Y}\|R_{X\mid Y}).\]
换句话说：

\textbf{联合 KL 散度总额已经确定。如果观察路径之后仍然保留了很多输入分布差异，那么路径本身所占用的 KL 散度预算就相应减少。}

接下来的问题变成：哪些输入信息会在观察完整路径后，仍然造成模型 $P$ 与 $R$ 之间的差异？

\subsection{为什么那些从未出现在小测试池中的物品特别重要？}
回顾我们之前定义的数量：
\[W=\mathbb E_P[w],\]
其中 $w$ 表示真实执行路径上从未出现在小测试池中的物品数量。

这些物品为什么会影响条件 KL 散度？

因为两个模型在完整路径终止后的状态不同。

在真实模型 $P$ 下，算法必须正确识别全部缺陷品。

因此，如果已经知道真实完整路径：
\[Y=y(D),\]
那么输入配置 $D$ 已经唯一确定：
\[P(X=D\mid Y=y(D))=1.\]
也就是说：
\[P_{X\mid Y=y(D)}\]
是集中在 $D$ 上的点质量分布。

但是，在人工模型 $R$ 下，大测试池的回答并不提供真实输入信息。

因此，即使完整路径已经确定，某些物品的状态仍然可能没有被确定。

特别是那些从未出现在小测试池中的物品。

由于它们在人工模型下仍具有随机性，因此：
\[R_{X\mid Y=y(D)}\]
不一定集中在 $D$ 上。

这就产生了条件 KL 散度。

我们需要将这种差异定量表示出来。

\subsection{怎样把未被小测试触及的物品数量转化为 KL 散度？}
固定真实输入 $D$ 及其完整路径 $y(D)$。

记：
\[V(D)=\text{这条路径上所有小规范化测试池的并集}.\]
定义：
\[w_1(D)=|D\setminus V(D)|.\]
它表示从未出现在小测试池中的缺陷品数量。

再定义：
\[w_0(D)=|([N]\setminus D)\setminus V(D)|.\]
它表示从未出现在小测试池中的正常品数量。

于是：
\[w(D)=w_1(D)+w_0(D).\]
现在我们考虑人工模型 $R$。

为了使输入配置与真实配置 $D$ 在所有小测试池上给出相同回答，一个充分条件是：它们在集合 $V(D)$ 中的全部坐标上完全一致。

但在 $V(D)$ 外，我们暂时不限制输入状态。

因此，在人工模型 $R$ 下，有一类输入能够产生路径 $y(D)$：它们在 $V(D)$ 内与 $D$ 完全相同，在外面任意取值，并且所有大测试池的公平硬币都给出指定回答。

这一类输入已经足以给出参考路径概率的下界：
\begin{equation}
\boxed{R_Y(y(D))\ge2^{-a(D)}p^{r-w_1(D)}q^{N-r-w_0(D)}.}
\tag{35}
\end{equation}
为什么是这个形式？

因为：
\begin{itemize}
\item $V(D)$ 内共有 $r-w_1(D)$ 件缺陷品，必须保持为缺陷品；
\item $V(D)$ 内共有 $N-r-w_0(D)$ 件正常品，必须保持为正常品；
\item $V(D)$ 外的状态不受限制，对所有可能取值求和后总概率为 1；
\item 所有大测试池的人工回答需要贡献 $2^{-a(D)}$。
\end{itemize}
而我们已经知道：
\[R_{X,Y}(D,y(D))=p^rq^{N-r}2^{-a(D)}.\]
因此：
\[
\begin{aligned}
R(X=D\mid Y=y(D))
&=\frac{R_{X,Y}(D,y(D))}{R_Y(y(D))}\\
&\le p^{w_1(D)}q^{w_0(D)}.
\end{aligned}
\]
也就是说：
\begin{equation}\boxed{R(X=D\mid Y=y(D))\le p^{w_1(D)}q^{w_0(D)}.}\tag{36}\end{equation}
现在使用点质量分布的 KL 散度公式。

由于真实条件分布集中在 $D$ 上：
\[
\begin{aligned}
&D_{\mathrm{KL}}(P_{X\mid Y=y(D)}\|R_{X\mid Y=y(D)})\\
&=-\log_2R(X=D\mid Y=y(D))\\
&\ge w_1(D)\log_2\frac1p+w_0(D)\log_2\frac1q.
\end{aligned}
\]
因此，每一件未被小测试池触及的物品，都能为条件 KL 散度贡献一定的下界。

定义：
\begin{equation}
\boxed{J=\mathbb E_P\left[w_1(X)\log_2\frac1p+w_0(X)\log_2\frac1q\right].}
\tag{37}
\end{equation}
那么：
\[\mathbb E_{Y\sim P_Y}D_{\mathrm{KL}}(P_{X\mid Y}\|R_{X\mid Y})\ge J.\]
代回式 (33)：
\begin{equation}\boxed{D_R+J\le K+A.}\tag{38}\end{equation}
这是人工参考模型比较中的核心结论。

它意味着：

\textbf{真实模型与人工模型之间的联合 KL 散度，总额为 $K+A$；其中至少有 $J$ 留在已知路径后的输入条件分布差异中，因此能够用于完整路径分布差异的部分至多为 $K+A-J$。}

这就是作者为什么不能简单地丢掉 KL 散度链式法则中的条件项。

保留它，反而能够得到更强的约束。

\subsection{怎样把 $J$ 与我们真正关心的 $W$ 联系起来？}
我们真正关心的是未被小测试池触及的物品数量：
\[W=\mathbb E_P[w_1(X)+w_0(X)].\]
而 $J$ 的定义中，每件未被触及的物品都带有一个对数权重。

由于：
\[0.249\le p\le0.251,\]
所以：
\[p\le0.251,\qquad q\le0.751.\]
论文通过精确数值比较得到：
\[\log_2\frac1p>\frac25,\qquad\log_2\frac1q>\frac25.\]
因此：
\[
\begin{aligned}
J&=\mathbb E_P\left[w_1(X)\log_2\frac1p+w_0(X)\log_2\frac1q\right]\\
&\ge\frac25\mathbb E_P[w_1(X)+w_0(X)].
\end{aligned}
\]
于是：
\begin{equation}\boxed{J\ge\frac25W.}\tag{39}\end{equation}
至此，我们终于能够通过 KL 散度控制未被小测试池触及的物品数量。

但还有一个任务：利用模型 $R$ 下全新物品对的不平衡性，控制 $B$。

\section{怎样把全新物品对 $B$ 与未触及物品 $W$ 一起控制？}
\subsection{再次使用加权二元不等式}
前面已经证明，对于任意全新物品对节点：
\[\rho_v^R=q^2,\]
并且：
\[\rho_v^R-\frac12>\frac3{50}.\]
因此，可以再次使用：
\[\omega(1-h(\pi_v))+d(\pi_v\|\rho_v^R)\ge\frac{2\omega}{(\omega+1)\ln2}\left(\rho_v^R-\frac12\right)^2.\]
这次论文选择：
\[\omega=64.\]
为什么与之前选择的 $1/10$ 不同？

因为我们现在要控制的是另一类测试池，而且后面需要将多个不等式组合起来。

权重 $\omega$ 影响信息缺口项与 KL 散度项的相对系数。

作者选择不同的权重，是为了让最终合并得到的线性系数足够好。

经过数值估计，可以得到：
\[\frac1{99}<64(1-h(\pi_v))+d(\pi_v\|\rho_v^R).\]
现在对所有全新物品对节点按照真实访问概率求和。

定义：
\[\Gamma_{\mathrm{fr}}=\sum_{v\text{ 为全新物品对节点}}P(v)(1-h(\pi_v)).\]
那么：
\[\frac B{99}\le64\Gamma_{\mathrm{fr}}+\sum_{v\text{ 为全新物品对节点}}P(v)d(\pi_v\|\rho_v^R).\]
接下来，我们需要控制最后这一项 KL 散度之和。

\subsection{为什么不能直接用整个 $D_R$？}
我们已经知道：
\[D_R=D_{\mathrm{KL}}(P_Y\|R_Y).\]
根据决策树 KL 散度链式法则：
\[D_R=\sum_vP(v)d(\pi_v\|\rho_v^R).\]
所以，全新物品对节点的 KL 散度之和肯定不超过 $D_R$。

直接使用：
\[\frac B{99}\le64\Gamma_{\mathrm{fr}}+D_R\]
当然也成立。

但作者没有停在这里。

因为还可以得到一个更强的估计。

回忆人工模型 $R$ 的构造：对于每个大测试池，它都使用公平硬币回答。

因此：
\[\rho_v^R=\frac12.\]
我们在第十课已经证明：
\[d\left(\pi_v\middle\|\frac12\right)=1-h(\pi_v).\]
因此，大测试池在 $D_R$ 中的全部贡献恰好为 $\Gamma_{\mathrm{lg}}$。

也就是说：
\[D_R=\Gamma_{\mathrm{lg}}+\text{其他节点的非负 KL 散度贡献}.\]
所以，全新物品对节点的 KL 散度之和最多为：
\[D_R-\Gamma_{\mathrm{lg}}.\]
这就得到：
\begin{equation}
\boxed{\frac B{99}\le64\Gamma_{\mathrm{fr}}+D_R-\Gamma_{\mathrm{lg}}.}
\tag{40}
\end{equation}
这里保留 $-\Gamma_{\mathrm{lg}}$ 非常重要。

因为它表示：大测试池已经占用了 $D_R$ 中确定的一部分 KL 散度，我们不能再把这部分预算全部用于全新物品对。

这一步体现了作者对信息论代价的精确分配。

\subsection{终于可以同时控制 $B$ 和 $W$}
此前，我们通过联合 KL 散度链式法则得到了：
\[D_R+J\le K+A.\]
而现在：
\[\frac B{99}\le64\Gamma_{\mathrm{fr}}+D_R-\Gamma_{\mathrm{lg}}.\]
因此：
\[\frac B{99}+J\le64\Gamma_{\mathrm{fr}}+K+A-\Gamma_{\mathrm{lg}}.\]
两边乘以 99：
\[B+99J\le6336\Gamma_{\mathrm{fr}}+99K+99A-99\Gamma_{\mathrm{lg}}.\]
由于：
\[J\ge\frac25W,\]
所以：
\[99J\ge\frac{198}{5}W\ge W.\]
最终得到：
\begin{equation}\boxed{B+W\le6336\Gamma_{\mathrm{fr}}+99K+99A-99\Gamma_{\mathrm{lg}}.}\tag{41}\end{equation}
现在，我们终于解决了控制小测试池的困难。

这个公式同时约束了：
\[B=\text{全新物品对平均次数},\]
以及：
\[W=\text{平均未触及物品数量}.\]
而右边只涉及：
\begin{itemize}
\item 真实算法的信息缺口；
\item 输入模型比较误差 $K$；
\item 已经得到约束的大测试池次数 $A$。
\end{itemize}
至此，第二条概率论主线的关键障碍全部解决。

接下来只需要将这些估计代回最初的组合计数不等式。

\section{为什么这些局部估计终于能够推出整棵树的下界？}
回忆我们最初从物品首次出现次数得到：
\[N-L\le B+W-A.\]
现在代入式 (41)：
\[
\begin{aligned}
N-L&\le B+W-A\\
&\le6336\Gamma_{\mathrm{fr}}+99K+98A-99\Gamma_{\mathrm{lg}}.
\end{aligned}
\]
这时，为什么之前需要控制大测试池的 $A$，就清楚了。

因为人工模型 $R$ 的联合 KL 散度会产生 $K+A$，导致 $B+W$ 的上界中出现正的 $99A$。

即使首次出现计数中有一个 $-A$，合并后仍然剩下 $98A$。

所以必须使用此前建立的大测试池上界：
\[A\le\frac{137}{2}\Gamma_{\mathrm{lg}}+685K.\]
代入：
\[
\begin{aligned}
N-L&\le6336\Gamma_{\mathrm{fr}}+99K\\
&\quad+98\left(\frac{137}{2}\Gamma_{\mathrm{lg}}+685K\right)-99\Gamma_{\mathrm{lg}}.
\end{aligned}
\]
整理系数：
\[98\cdot\frac{137}{2}-99=6614,\]
\[99+98\cdot685=67229.\]
所以：
\begin{equation}\boxed{N-L\le6336\Gamma_{\mathrm{fr}}+6614\Gamma_{\mathrm{lg}}+67229K.}\tag{42}\end{equation}
现在回忆：
\[\Gamma=\sum_vP(v)(1-h(\pi_v)).\]
大测试池与全新物品对是两类互不重叠的测试节点。

因此：
\[\Gamma_{\mathrm{lg}}+\Gamma_{\mathrm{fr}}\le\Gamma.\]
又因为：
\[6336<6614,\]
所以：
\[6336\Gamma_{\mathrm{fr}}+6614\Gamma_{\mathrm{lg}}\le6614\Gamma.\]
最终：
\begin{equation}\boxed{N-L\le6614\Gamma+67229K.}\tag{43}\end{equation}
而：
\[\Gamma=L-H.\]
所以：
\[N-H=(N-L)+(L-H)=(N-L)+\Gamma.\]
于是：
\[N-H\le6615\Gamma+67229K.\]
整理得到：
\begin{equation}\boxed{6615(L-H)\ge N-H-67229K.}\tag{44}\end{equation}
这就是论文定理 9.1 的核心不等式。

它适用于缺陷密度位于：
\[0.249\le p\le0.251\]
的所有正确有限确定性自适应分组测试决策树。

\subsection{现在回顾为什么需要这些工具}
这里我们可以完整回顾一次概率论主线。

我们最初希望证明：
\[L-H=\Gamma\]
不可能太小。

但真实模型 $P$ 下的条件测试概率难以计算。

于是，我们先观察到一个纯组合关系：
\[N-L\le B+W-A.\]
这使我们把问题转化为控制大测试池、全新物品对和未触及物品的数量。

接下来：

为了控制大测试池，我们引入独立模型 $Q$。

但 $Q$ 与真实模型不同，于是通过 KL 散度控制累计模型差异。

要将比较模型的不平衡性转化为真实信息缺口，就引入加权二元不等式。

而自适应历史影响条件独立性，于是又使用单调事件关联不等式。

这样，终于控制了 $A$。

随后发现，$Q$ 无法直接控制全新物品对，于是设计人工模型 $R$。

为了证明 $R$ 下的全新物品对仍然具有独立伯努利条件分布，作者建立固定历史条件分布公式。

但 $R$ 不是真实测试模型，不能直接使用原来的 $K$ 上界，于是计算联合 KL 散度 $K+A$。

为了进一步控制 $W$，作者保留 KL 链式法则中的条件输入散度，得到 $J\ge(2/5)W$。

最后，再利用大测试池在人工模型中恰好贡献信息缺口的事实，控制 $B$ 并得到联合上界。

最终通过首次出现计数，得到整棵决策树的信息论不等式。

\textbf{整条路线中的每一次新工具引入，都是为了解决上一步尚未解决的具体问题。}

现在，信息论部分已经完成。

但我们的最终证明还差最后一步。

我们必须把式 (44) 转化为节省量上界，再与结构递推强迫的下界比较。

\section{回到最初的反证：信息论究竟允许节省多少次测试？}
\subsection{我们现在已经获得了什么？}
首先回顾两条主线分别得到了什么。

\textbf{第一条：结构递推。}

从假设：
\[\Delta(d,n)\ge1\]
出发，我们证明存在一个新的问题 $(r,N)$，满足：
\[N>2^{29},\qquad0.249<\frac rN\le0.25,\]
并且：
\begin{equation}\boxed{\Delta(r,N)\ge1+g_\star N.}\tag{45}\end{equation}
其中：
\[g_\star=\frac{7441}{39434}.\]
这是节省量的下界。

\textbf{第二条：信息论分析。}

对于密度接近 $1/4$ 的问题，我们证明：
\begin{equation}\boxed{6615(L-H)\ge N-H-67229K.}\tag{46}\end{equation}
其中：
\[H=\log_2\binom Nr,\]
\[K=Nh(p)-H,\qquad p=\frac rN.\]
现在的任务非常明确：

\textbf{从式 (46) 推导出节省量 $\Delta(r,N)$ 的上界，再证明它与式 (45) 不可能同时成立。}

\subsection{为什么平均测试次数下界可以限制最坏情况节省量？}
由式 (46)：
\[L-H\ge\frac{N-H-67229K}{6615}.\]
所以：
\[L\ge H+\frac{N-H-67229K}{6615}.\]
而：
\[M(r,N)\ge L.\]
因此：
\begin{equation}\boxed{M(r,N)\ge H+\frac{N-H-67229K}{6615}.}\tag{47}\end{equation}
我们真正关心的是：
\[\Delta(r,N)=N-1-M(r,N).\]
所以：
\[
\begin{aligned}
\Delta(r,N)&\le N-1-H-\frac{N-H-67229K}{6615}\\
&=\frac{6614}{6615}(N-H)+\frac{67229}{6615}K-1.
\end{aligned}
\]
又因为：
\[H=Nh(p)-K,\]
所以：
\[N-H=N(1-h(p))+K.\]
代入：
\[
\begin{aligned}
\Delta(r,N)&\le\frac{6614}{6615}N(1-h(p))\\
&\quad+\frac{6614+67229}{6615}K-1.
\end{aligned}
\]
为了简化记号，定义：
\begin{equation}\boxed{\alpha=\frac{6614}{6615},\qquad\beta=\frac{73843}{6615}.}\tag{48}\end{equation}
其中：
\[\alpha<1,\qquad\beta<12.\]
于是：
\[\Delta(r,N)\le\alpha N(1-h(p))+\beta K-1.\]
而：
\[K\le\log_2(N+1).\]
最终：
\begin{equation}\boxed{\Delta(r,N)\le\alpha N(1-h(p))+\beta\log_2(N+1)-1.}\tag{49}\end{equation}
这就是我们想要的信息论节省量上界。

\subsection{为什么这里的 $\alpha<1$ 如此关键？}
回顾最初的普通熵下界。

它只能给出近似形式：
\[\Delta(r,N)\lesssim N(1-h(p)).\]
现在我们得到：
\[\Delta(r,N)\le\alpha N(1-h(p))+O(\log N),\]
其中：
\[\alpha=\frac{6614}{6615}<1.\]
也就是说，经过前面的全部信息论分析，我们把最大允许节省量的主要线性项系数，从 1 改善成了一个严格小于 1 的数。

虽然：
\[1-\alpha=\frac1{6615}\]
非常小，但这正是整篇论文需要的信息论改进。

因为结构递推的节省量下界系数 $g_\star$，与普通熵允许的节省量系数 $1-h(1/4)$ 本来就十分接近。

只要这个微小改进能够使两个系数的大小关系发生逆转，我们就有希望得到矛盾。

不过，在正式比较之前，还要解决一个小问题。

\subsection{为什么要把实际停止密度 $p$ 换成 $1/4$？}
我们知道：
\[0.249<p\le\frac14.\]
但节省量上界包含：
\[1-h(p),\]
而我们希望与固定的常数 $g_\star$ 比较。

如果直接使用不同的 $p$，比较会不方便。

因此，我们希望将 $1-h(p)$ 统一替换成 $1-h(1/4)$，同时控制替换产生的误差。

这就是前面整数舍入控制在最后一步再次发挥作用的原因。

回忆：
\[\frac14-\frac3{4N}<p\le\frac14.\]
我们又知道：
\[h'(u)=\log_2\frac{1-u}{u}.\]
在：
\[0.249\le u\le\frac14\]
上，有：
\[0<h'(u)<2.\]
因此：
\[h(1/4)-h(p)\le2\left(\frac14-p\right).\]
而：
\[\frac14-p<\frac3{4N}.\]
所以：
\[h(1/4)-h(p)<\frac3{2N}.\]
整理：
\[1-h(p)<1-h(1/4)+\frac3{2N}.\]
两边乘以 $N$：
\begin{equation}\boxed{N(1-h(p))<N(1-h(1/4))+\frac32.}\tag{50}\end{equation}
这个结果意味着：

虽然停止密度不一定恰好等于 $1/4$，但在节省量上界中，把它替换为 $1/4$，最多只会产生一个固定大小的误差。

这个误差不会随着 $N$ 线性增长。

因此，它不会破坏我们最后需要的线性项比较。

\subsection{得到可以直接比较的两个节省量界}
定义：
\begin{equation}\boxed{e_{1/4}=1-h(1/4).}\tag{51}\end{equation}
这是一个固定常数。

我们已经知道：
\[e_{1/4}\approx0.1887218755.\]
结合式 (49) 和 (50)：
\begin{equation}
\boxed{\Delta(r,N)<\alpha e_{1/4}N+\frac{3\alpha}{2}+\beta\log_2(N+1)-1.}
\tag{52}
\end{equation}
另一方面，结构递推要求：
\begin{equation}\boxed{\Delta(r,N)\ge1+g_\star N.}\tag{53}\end{equation}
现在，我们终于可以把两条证明主线直接放在一起比较。

\section{最后的矛盾究竟出现在哪里？}
\subsection{两个节省量约束必须同时成立}
如果原始反证假设成立，那么经过结构递推得到的新问题 $(r,N)$，必须同时满足：
\[1+g_\star N\le\Delta(r,N)\]
和：
\[\Delta(r,N)<\alpha e_{1/4}N+\frac{3\alpha}{2}+\beta\log_2(N+1)-1.\]
因此：
\[1+g_\star N<\alpha e_{1/4}N+\frac{3\alpha}{2}+\beta\log_2(N+1)-1.\]
将线性项移到左边：
\begin{equation}\boxed{(g_\star-\alpha e_{1/4})N<\beta\log_2(N+1)+\frac{3\alpha}{2}-2.}\tag{54}\end{equation}
由于：
\[0<\alpha<1,\]
所以：
\[\frac{3\alpha}{2}-2<0.\]
因此，如果原始反证假设成立，那么必然有：
\begin{equation}\boxed{(g_\star-\alpha e_{1/4})N<\beta\log_2(N+1).}\tag{55}\end{equation}
现在，所有复杂的分组测试结构、概率模型和信息论工具，都已经被压缩成了这个简单的数值不等式。

我们只需要证明它不可能成立。

\subsection{为什么一定要验证 $g_\star-\alpha e_{1/4}>0$？}
观察式 (55)：
\[(g_\star-\alpha e_{1/4})N<\beta\log_2(N+1).\]
如果：
\[g_\star-\alpha e_{1/4}\le0,\]
那么左边不是一个严格正的线性项，就无法利用线性函数增长快于对数函数的事实完成反证。

因此，我们必须证明：
\[\boxed{g_\star-\alpha e_{1/4}>0.}\]
这就是为什么论文最后需要非常精确的数值计算。

我们已经知道：
\[g_\star=\frac{7441}{39434}\approx0.1886950347,\]
\[e_{1/4}=1-h(1/4)\approx0.1887218755,\]
\[\alpha=\frac{6614}{6615}\approx0.9998488284.\]
因此：
\[\alpha e_{1/4}\approx0.1886933462.\]
所以：
\[g_\star-\alpha e_{1/4}\approx1.68858\times10^{-6}.\]
这个差值虽然非常小，但确实为正。

论文为了避免依靠近似小数证明严格不等式，在附录中使用有限对数级数及其余项界，给出精确的有理数证书，证明：
\begin{equation}\boxed{g_\star-\alpha e_{1/4}>10^{-6}.}\tag{56}\end{equation}
这一步非常关键。

它不仅说明差值为正，还给出一个明确的正下界。

\subsection{为什么还需要规模条件 $N>2^{29}$？}
即使已经证明：
\[g_\star-\alpha e_{1/4}>10^{-6},\]
式 (55) 也只是要求：
\[\frac{N}{10^6}<\beta\log_2(N+1).\]
现在我们需要证明：当 $N$ 足够大时，这个不等式不可能成立。

因为：
\[\beta<12,\]
所以只需要证明：
\begin{equation}\boxed{12\log_2(N+1)<\frac{N}{10^6}.}\tag{57}\end{equation}
论文证明，这个不等式对于：
\[N\ge2^{29}\]
成立。

为什么？

我们可以先看最小规模：
\[N=2^{29}.\]
由于：
\[2^{29}=536870912,\]
所以：
\[\frac N{10^6}>536.\]
另一方面：
\[2^{29}+1<2^{30}.\]
因此：
\[12\log_2(N+1)<12\times30=360.\]
所以在 $N=2^{29}$ 时：
\[12\log_2(N+1)<\frac N{10^6}.\]
而对于更大的 $N$，左边只有对数增长，右边是线性增长。

论文还通过导数严格证明：两边的差值在该范围内持续增加。

因此，式 (57) 对全部 $N\ge2^{29}$ 成立。

这里就解释了为什么前面必须花费篇幅证明递推后的规模仍然满足：
\[N>2^{29}.\]
如果没有这个规模保证，就不能直接使用最后的数值不等式完成反证。

\subsection{最终矛盾}
现在我们已经得到：
\[g_\star-\alpha e_{1/4}>10^{-6},\]
\[N>2^{29},\]
以及：
\[\beta<12.\]
因此：
\[
\begin{aligned}
(g_\star-\alpha e_{1/4})N&>\frac N{10^6}\\
&>12\log_2(N+1)\\
&>\beta\log_2(N+1).
\end{aligned}
\]
即：
\begin{equation}\boxed{(g_\star-\alpha e_{1/4})N>\beta\log_2(N+1).}\tag{58}\end{equation}
但是，从原始反证假设出发，我们又推导出了：
\[\boxed{(g_\star-\alpha e_{1/4})N<\beta\log_2(N+1).}\]
两个不等式直接矛盾。

因此，最初的假设：
\[\Delta(d,n)\ge1\]
不可能成立。

由于 $\Delta(d,n)$ 是非负整数，所以只能有：
\[\Delta(d,n)=0.\]
根据定义：
\[\Delta(d,n)=n-1-M(d,n),\]
因此：
\[\boxed{M(d,n)=n-1.}\]
最终定理得证。

\section{现在重新审视整篇论文：各个工具为什么不可缺少？}
至此，我们已经从最初的反证假设一步步走到了最终结论。

现在回顾整条推导，就能够更清楚地区分：

哪些工具是为了研究分组测试的组合结构，哪些工具是为了进行概率比较，哪些工具用于把两种方法组合起来。

\begin{proofbox}\centering
\textbf{最终证明目标}
\[M(d,n)=n-1\]
\end{proofbox}
\[\Downarrow\]
\begin{proofbox}
\textbf{困难一：直接分析高密度问题的所有算法很困难。}

引入两特殊标签恒等式及结构递推，转向更小的问题。
\end{proofbox}
\[\Downarrow\]
\begin{proofbox}
\textbf{困难二：递推可能提前停止，或者不能准确到达目标密度。}

使用边界分析、规模估计和整数舍入控制。
\end{proofbox}
\[\Downarrow\]
\begin{proofbox}\centering
\textbf{得到密度约 $1/4$ 的新问题}
\[\Delta(r,N)\ge1+g_\star N\]
\end{proofbox}
\[\Downarrow\]
\begin{proofbox}
\textbf{困难三：普通熵下界不足以排除这个节省量。}

研究额外信息缺口 $\Gamma=L-H$。
\end{proofbox}
\[\Downarrow\]
\begin{proofbox}
\textbf{困难四：信息缺口很难直接与总测试次数建立足够强的联系。}

对测试池分类，利用物品首次出现次数，得到 $N-L\le B+W-A$。
\end{proofbox}
\[\Downarrow\]
\begin{proofbox}
\textbf{困难五：真实模型下的大测试池条件概率很难计算。}

引入独立模型 $Q$、KL 散度比较及单调事件关联不等式，控制 $A$。
\end{proofbox}
\[\Downarrow\]
\begin{proofbox}
\textbf{困难六：真实独立模型无法直接控制全新物品对。}

构造人工模型 $R$，让大测试池使用公平回答，保留相关物品的条件独立性。
\end{proofbox}
\[\Downarrow\]
\begin{proofbox}
\textbf{困难七：无法直接用原来的 $K$ 控制人工模型的路径差异。}

计算联合 KL 散度 $K+A$，并保留条件输入分布差异，控制 $B+W$。
\end{proofbox}
\[\Downarrow\]
\begin{proofbox}\centering
\textbf{得到全树信息不等式}
\[6615(L-H)\ge N-H-67229K\]
\end{proofbox}
\[\Downarrow\]
\begin{proofbox}
\textbf{最后困难：信息论上界是否真的强于递推下界？}

利用严格数值证书，证明正的线性差距压倒对数误差。
\end{proofbox}
\[\Downarrow\]
\begin{proofbox}\centering
\textbf{产生矛盾，反证完成}
\[\boxed{M(d,n)=n-1}\]
\end{proofbox}

从这条逻辑链可以看出，整篇论文中最重要的并不是某个单独的不等式，而是几个相互配合的分析思想。

\subsection{结构递推负责制造一个强制要求}
我们最初只是知道：如果原始定理不成立，那么某个高密度问题可以节省至少一次测试。

这本身并不是很强的假设。

通过结构递推，我们把这个假设放大成：
\[\Delta(r,N)\ge1+g_\star N.\]
这样，原本只要求节省一次测试的假设，被转化为一个新问题中必须节省线性数量测试的要求。

这使反证变得有可能。

\subsection{信息论负责证明这个要求无法实现}
信息论部分没有直接计算最优测试算法，也没有枚举决策树。

它研究的是：所有合法的决策树都必须遵守怎样的概率和信息约束？

特别是通过：
\[\Gamma=L-H,\]
将问题转化为证明真实算法不可能避免一定的信息缺口。

最终得到：
\[6615\Gamma\ge N-H-67229K.\]
这限制了算法可能节省的测试次数。

\subsection{两个辅助模型分别解决不同困难}
独立模型 $Q$ 的作用是：在适当处理自适应历史后，控制大测试池的阴性概率。

人工模型 $R$ 的作用是：消除大测试池回答对输入产生的条件约束，使全新物品对能够保留独立性，并且使未被小测试触及的物品能够贡献条件 KL 散度。

两者的功能不同。

如果只使用 $Q$，难以处理全新物品对。

如果只使用 $R$，又不能直接得到大测试池在真实独立 OR 模型下的概率约束。

因此，作者把不同的模型用于不同类型的测试池。

\subsection{KL 散度并不是直接代表测试次数}
这点尤其值得强调。

在这篇论文中，我们从来没有直接写：
\[\text{KL 散度}=\text{需要增加的测试次数}\]
这样的结论。

KL 散度的作用是控制不同模型在整个测试过程中的累计概率差异。

通过加权二元不等式，这种概率差异与真实模型的信息缺口联系起来。

再通过物品首次出现次数的组合不等式，信息缺口才最终与测试次数建立联系。

因此，真正的逻辑是：
\[
\boxed{\text{模型差异}\longrightarrow\text{信息缺口约束}\longrightarrow\text{测试结构约束}\longrightarrow\text{测试次数下界}.}
\]
而不是简单地把 KL 散度解释成测试次数。

\section{最终总结：这篇论文的证明思想是什么？}
现在，我们可以不依赖具体章节顺序，用一段完整的话概括整篇证明。

论文首先假设存在一种能够在最坏情况下优于逐个测试的自适应分组测试算法。为了避免直接分析复杂的高密度决策树，作者利用分组测试的合法 OR 结构，建立不同规模问题之间的递推关系。通过这个递推，可以将原始高密度问题转化为一个缺陷密度约为 $1/4$、规模仍然足够大的问题，并且迫使新问题具有线性阶的测试节省量。

接下来，作者试图利用信息论排除这种节省量。但普通熵下界还不够强，因此转而研究真实算法超过熵下界的额外信息缺口。为了把信息缺口与测试次数联系起来，作者按照规范化后的测试池大小及物品首次出现情况进行分类，建立了涉及大测试池、全新物品对和未触及物品数量的组合计数关系。

随后，为了克服真实精确缺陷数量模型中的概率依赖困难，作者引入独立伯努利模型，并利用 KL 散度、单调事件关联不等式和加权二元不等式控制大测试池。对于独立模型仍然难以处理的全新物品对，作者进一步构造人工参考实验，让大测试池使用公平硬币回答，从而保留相关物品的条件独立性。再通过联合 KL 散度的精确计算及链式分解，控制全新物品对和未触及物品的总贡献。

这些局部估计最终与物品首次出现次数的不等式结合，得到一个适用于整棵决策树的信息论下界。该下界虽然只对密度接近 $1/4$ 的问题成立，但它恰好足以限制新问题可能节省的测试次数。

最后，通过精确的常数比较，作者证明：结构递推强迫的节省量，严格大于信息论允许的最大节省量。当问题规模足够大时，这个正的线性差距必然压倒剩余的对数误差，从而产生矛盾。

因此，最初假设存在更优算法是不可能的，最终证明：
\[
\boxed{M(d,n)=n-1\qquad\text{当 }n\ge2^{31},\ 0<d<n,\ \frac dn\ge0.369056.}
\]
这里也需要明确论文的证明范围：该结论尚未覆盖 Hu--Hwang--Wang 猜想所关注的完整密度区间 $d/n\ge1/3$，而是证明了一个严格高于 $1/3$ 的充分密度阈值。

\textbf{整篇论文最值得学习的核心方法是：先通过组合递推放大反证假设所要求的节省量，再通过精心设计的概率比较实验，证明任何合法自适应决策树都无法达到这一要求。}

从这个角度看，熵、KL 散度、条件独立性、人工参考实验和结构递推并不是彼此孤立的工具，而是分别解决整条证明路线中不同的障碍，最终共同完成了下界证明。

\end{document}
